\addchap{Notation}

One notation is used throughout the book. Where a chapter needs a local symbol that clashes with this list, the chapter says so; where a source paper uses a different notation, the chapter's notes translate it.

\newcommand{\notationrow}[2]{#1 & #2\\}
\newcommand{\notationhead}[1]{\multicolumn{2}{@{}l}{\textsc{\color{sea}#1}}\\[0.3ex]}
\noindent{\small
\renewcommand{\arraystretch}{1.1}%
\begin{tabular}{@{}>{\raggedright\arraybackslash}p{32mm}>{\raggedright\arraybackslash}p{84mm}@{}}
\notationhead{sizes}
\notationrow{$d$}{number of parameters, the dimension of the optimization variable}
\notationrow{$W\in\R^{d_{\mathrm{out}}\times d_{\mathrm{in}}}$}{a weight matrix; $G\in\R^{m\times n}$ a generic matrix or matrix gradient}
\notationrow{$N$, $l$}{number of layers and the index of a layer; widths $d_0,\dots,d_N$}
\notationrow{$n$, $B$, $K$}{number of samples, batch size, number of classes}
\notationrow{$T$, $t$}{number of steps and the step, counted from $t=0$; $t$ is continuous for flows; the online-learning arguments of \cref{ch:adaptive,ch:kronecker,ch:schedules} count rounds from $t=1$, as their sources do}
\addlinespace
\notationhead{objective}
\notationrow{$\Loss(\theta)$}{the objective: the training or the population loss, as each chapter says}
\notationrow{$\ell(\theta;\xi)$}{the loss on one example or one batch $\xi$}
\notationrow{$\Loss^\star$, $\theta^\star$}{the infimum of $\Loss$ and a minimizer; a superscript star marks an optimal value}
\notationrow{$\Delta_t$}{the suboptimality $\Loss(\theta_t)-\Loss^\star$}
\notationrow{$\theta_\star$}{a subscript star marks a target or a teacher}
\addlinespace
\notationhead{gradients and curvature}
\notationrow{$g_t$, $\zeta_t$}{a stochastic gradient at $\theta_t$, and its noise $g_t-\nabla\Loss(\theta_t)$}
\notationrow{$\Sigma(\theta)$}{the per-example covariance of the gradient, so that $\Cov(g_t)=\Sigma/B$}
\notationrow{$H$, $S$}{the Hessian $\nabla^2\Loss$ and the sharpness $S=\lmax(H)$}
\notationrow{$S_P$}{the preconditioned sharpness $\lmax(P^{-1}H)$}
\addlinespace
\notationhead{optimizers}
\notationrow{$\eta$, $\eta_t$}{step size (learning rate) and its schedule; $\eta_{\max}$ the peak, $T_{\mathrm w}$ the warmup length}
\notationrow{$\beta$; $\beta_1$, $\beta_2$, $\epsilon$}{momentum; Adam's averaging coefficients and its $\epsilon$}
\notationrow{$m_t$, $v_t$; $\hat m_t$, $\hat v_t$}{Adam's moving averages of $g$ and $g^2$, and their bias-corrected versions}
\notationrow{$\lambda$}{decoupled weight decay, $\theta\leftarrow(1-\eta\lambda)\theta-\eta u$; eigenvalues always carry a subscript, as in $\lambda_i$, $\lmax$, $\lmin$}
\end{tabular}

\medskip\noindent
\begin{tabular}{@{}>{\raggedright\arraybackslash}p{32mm}>{\raggedright\arraybackslash}p{84mm}@{}}
\notationhead{norms}
\notationrow{$\norm{\cdot}$, $\norm{\cdot}_*$}{a generic norm and its dual norm; an unadorned norm is never silently Euclidean}
\notationrow{$\norm{\cdot}_2$, $\norm{\cdot}_1$, $\norm{\cdot}_\infty$}{the $\ell_2$, $\ell_1$ and $\ell_\infty$ norms of a vector}
\notationrow{$\norm{W}_{\op}$, $\norm{W}_F$, $\norm{W}_{\nuc}$}{spectral, Frobenius and nuclear norms of a matrix; $\norm{W}_{p\to q}$ an induced norm}
\notationrow{$\norm{x}_{\rms}$}{the root-mean-square norm $\norm{x}_2/\sqrt{\dim x}$; $\norm{W}_{\rms\to\rms}=\sqrt{d_{\mathrm{in}}/d_{\mathrm{out}}}\,\norm{W}_{\op}$}
\notationrow{$\norm{x}_P$}{$\sqrt{x^\top Px}$ for a positive definite $P$}
\addlinespace
\notationhead{steepest descent}
\notationrow{$\lmo(g)$}{a minimizer of $\ip{g,u}$ over $\norm{u}\le1$, so that $\ip{g,\lmo(g)}=-\norm{g}_*$}
\notationrow{$g^\sharp$}{the sharp operator $\argmax_u\{\ip{g,u}-\tfrac12\norm{u}^2\}=-\norm{g}_*\lmo(g)$}
\notationrow{normalized step}{$\theta_{t+1}=\theta_t+\eta\lmo(g_t)$: a step of length $\eta$ in the chosen norm}
\notationrow{unnormalized step}{$\theta_{t+1}=\theta_t-\eta g_t^\sharp$}
\notationrow{$\sign$, $\polar$}{coordinatewise sign with $\sign0=0$; $\polar(G)=UV^\top$ for a compact singular value decomposition $G=U\diag(\sigma)V^\top$}
\addlinespace
\notationhead{smoothness}
\notationrow{$L$-smooth}{$\norm{\nabla\Loss(\theta)-\nabla\Loss(\theta')}_*\le L\norm{\theta-\theta'}$, with respect to a stated norm}
\notationrow{$(L_0,L_1)$}{$\norm{\nabla^2\Loss(\theta)}\le L_0+L_1\norm{\nabla\Loss(\theta)}$ (\cref{ch:spikes})}
\notationrow{$(K_0,K_1)$}{$\norm{\nabla^2\Loss(\theta)}\le K_0+K_1(\Loss(\theta)-\Loss^\star)$ (\cref{ch:warmup})}
\addlinespace
\notationhead{matrices, probability, asymptotics}
\notationrow{$\otimes$, $\vect$, $\preceq$}{Kronecker product, column-stacking vectorization, Löwner order}
\notationrow{$\cN$, $\E$, $\Var$, $\Cov$}{Gaussian law, expectation, variance, covariance}
\notationrow{$a\lesssim b$, $a\asymp b$}{$a\le Cb$, and $a\lesssim b\lesssim a$, for a constant $C$ depending only on parameters declared fixed}
\addlinespace
\notationhead{status of an argument}
\notationrow{$\blacksquare$}{end of a complete proof}
\notationrow{proof sketch}{an argument that omits steps, with a pointer to where they are}
\notationrow{$\Diamond$}{end of a heuristic derivation: an argument that is not a proof}
\notationrow{observation}{an empirical regularity, with its evidence}
\notationrow{folklore}{a rule used in practice for which no evidence can be cited}
\end{tabular}\par}
